\section{코드 별 연계 구조}

본 장은 앞 장들에서 개별 분석한 구성요소들이 실제 동작 시 어떻게 맞물리는지를
종단 간(end-to-end) 연계 관점에서 정리한다.

\subsection{종단 간 명령 처리 시퀀스}

블록 하나가 실행되어 하드웨어가 반응하기까지의 호출 사슬은 다음과 같다.

\captionof{listing}{블록 1개 실행의 호출 사슬}
\begin{minted}{text}
[웹]
  main.js : 실행 버튼
    -> commandexecutor.processBlock(block)
       -> generateCommand(block)        // "SERVO_tFORWARD,2"
       -> sendCommand(cmd)              // fire-and-forget 여부 판단
          -> bluetooth.sendData(cmd, wait)
             -> 20바이트 청크 write (0xFFE1 특성)
[BLE 모듈] -> UART 9600
[피코]
  main.py : _read_uart_line()           // 개행까지 버퍼 조립
    -> CommandProcessor.process("SERVO_tFORWARD,2")
       -> _handle_timed_wheel()
          -> robot.wheel.forward(speed)  // 2초 차단 후 정지
    -> 응답 필요 시 _send_response("1")  // 20바이트 청크
[웹]
  bluetooth.handleRxData() -> pendingResolve("1")  // 프라미스 해소
\end{minted}

같은 사슬을 시간 축 시퀀스로 표현하면 그림~\ref{fig:e2e}와 같다. 6개 참여
요소가 호출(실선)과 응답(점선)으로 이어지며, 시간 명령(\code{tFORWARD,2})의
경우 \code{wheel.forward}가 2초간 펌웨어를 점유한 뒤에야 응답이 역방향으로
전파되는 것을 알 수 있다.

\begin{diagram}{종단 간 명령\textbf{·}응답 시퀀스 연계도}{fig:e2e}
\node[dom, text width=22mm] (ce) at (0,0) {command-\\executor};
\node[dom, text width=18mm] (bt) at (3.0,0) {bluetooth};
\node[boxw, text width=14mm] (ble) at (6.0,0) {BLE 모듈};
\node[dom, text width=16mm] (mp) at (9.0,0) {main.py};
\node[dom, text width=20mm] (cp) at (12.2,0) {Command-\\Processor};
\node[hw, text width=16mm] (hw) at (15.2,0) {하드웨어};
\foreach \n in {ce,bt,ble,mp,cp,hw} \draw[rulegray,dashed] (\n.south) -- ++(0,-7.2);
\newcommand{\seq}[5]{\draw[#5] ([yshift=#3cm]#1.south) -- node[lbl,above]{\scriptsize #4} ([yshift=#3cm]#2.south);}
\seq{ce}{bt}{-0.8}{sendData}{flow}
\seq{bt}{ble}{-1.5}{20B write}{flow}
\seq{ble}{mp}{-2.2}{UART 9600}{flow}
\seq{mp}{cp}{-2.9}{process()}{flow}
\seq{cp}{hw}{-3.6}{wheel.forward (2초)}{flow}
\seq{hw}{cp}{-4.3}{완료}{back}
\seq{cp}{mp}{-5.0}{"1"}{back}
\seq{mp}{ble}{-5.7}{20B}{back}
\seq{ble}{bt}{-6.4}{notify}{back}
\seq{bt}{ce}{-7.1}{pendingResolve}{back}
\end{diagram}

\subsection{모듈 간 결합 지점}

\begin{longtable}{L{4.6cm} L{4.4cm} L{4.4cm}}
\toprule
\rowcolor{tablehead}\textbf{결합 지점} & \textbf{웹 측} & \textbf{펌웨어 측} \\
\midrule
\endhead
명령 프로토콜 문자열 & \code{commandexecutor.js} & \code{process\_data.py} \\
BLE UUID/청크 규약 & \code{constants.js}, \code{bluetooth.js} & \code{main.py}(UART) \\
센서 변수 반영 & \code{bluetooth.js} $\to$ \code{state.variables} & \code{DIST:}/\code{MAG:} 응답 \\
시스템 설정 동기화 & \code{dashboard.html} & \code{GET\_SYS}/\code{SYS\_SET} \\
모듈 활성 상태 & 대시보드 토글 & \code{system.json}/\code{SET\_MODULE} \\
\bottomrule
\end{longtable}

\subsection{명령 프로토콜: 두 코드베이스의 계약}

웹과 펌웨어는 코드를 공유하지 않으므로, \emph{명령 문자열 프로토콜}이 둘
사이의 유일한 계약(contract)이다. 따라서 한쪽의 명령 형식 변경은 반드시
다른 쪽과 함께 갱신되어야 한다.

\begin{infobox}[title=프로토콜 변경 시 동기화 규칙]
\begin{itemize}
  \item 웹에서 새 명령을 생성하면(\code{commandexecutor.js}),
        펌웨어 핸들러(\code{process\_data.py})를 함께 추가한다.
  \item 응답을 기대하는 명령은 양측의 fire-and-forget 목록
        (\code{FIRE\_AND\_FORGET\_HEADS} $\leftrightarrow$ \code{NO\_RESPONSE\_CMDS})을
        일치시켜 명령-응답 오정렬을 방지한다.
  \item 인자 구분자(쉼표)\textbf{·}종단(개행)\textbf{·}20바이트 청크 규약을 양측이
        동일하게 따른다.
  \item 프로토콜 명세는 \code{Document/API\_DOCUMENTATION.md}를 단일 출처로 유지한다.
\end{itemize}
\end{infobox}

\subsection{상태 공유와 데이터 결합 (웹 내부)}

웹 클라이언트 내부에서는 \code{state.js}가 사실상의 공유 데이터 허브이다.
\code{bluetooth.js}가 수신한 센서 값을 \code{state.variables}에 기록하면,
\code{commandexecutor.js}의 조건식 평가가 이를 읽어 제어 흐름을 결정한다.
대시보드(iframe)와 본체는 \code{postMessage}로 \code{STATUS} 텔레메트리를
주고받아 상태를 동기화한다.

\subsection{설정의 이중 소유와 정합성}

시스템 설정은 두 곳에 존재한다. 펌웨어의 \code{system.json}(정본)과 웹의
\code{localStorage}(\code{ares-system-config}, 캐시\textbf{·}UI 기본값)이다.
대시보드는 \code{GET\_SYS}로 펌웨어 값을 읽어 표시하고, \code{SYS\_SET}/
\code{CALIB\_SET}으로 펌웨어에 반영한다. 따라서 펌웨어 값이 실효(즉 실제
적용되는) 기준이며, 웹은 이를 조회\textbf{·}갱신하는 클라이언트로 동작한다(그림~\ref{fig:cfgsync}).

\begin{diagram}{설정의 이중 소유와 동기화 연계도}{fig:cfgsync}
\node[store, text width=26mm] (json) {system.json\\{\scriptsize 펌웨어 정본}};
\node[dom, right=28mm of json, text width=26mm] (dash) {대시보드\\{\scriptsize dashboard.html}};
\node[store, right=24mm of dash, text width=26mm] (ls) {localStorage\\{\scriptsize ares-system-config}};
\draw[flow] (json.north east) to[bend left=18] node[lbl,above]{GET\_SYS (조회)} (dash.north west);
\draw[back] (dash.south west) to[bend left=18] node[lbl,below]{SYS\_SET / CALIB\_SET (반영)} (json.south east);
\draw[flow] (dash) -- node[lbl,above]{UI 기본값 캐시} (ls);
\node[font=\scriptsize\sffamily\color{accentcopper}, below=2mm of json] {실효 기준};
\end{diagram}

\subsection{시뮬레이션과 실기의 분기}

\code{commandexecutor.js}의 디스패치는 시뮬레이션 여부에 따라 갈린다.
시뮬레이션 모드에서는 동일한 명령 문자열이 BLE 대신 \code{simulation.js}의
싱크로 전달되어 3D로 시연된다. 즉, \emph{명령 생성 로직은 단일}이고
\emph{전송 대상만 교체}되므로, 시뮬레이터와 실기의 동작 의미가 일치한다.
이는 ``동일 블록 $\to$ 동일 명령 $\to$ (시뮬 또는 실기)'' 구조로, 교육 현장에서
로버 없이도 동일한 코딩 경험을 제공하는 핵심 설계이다.

\begin{teachbox}{두 코드베이스를 잇는 계약}
\teachhead{설계 의도}
\begin{itemize}
  \item 웹과 펌웨어의 \emph{유일한 계약}은 명령 문자열이다. 한쪽을 바꾸면 반드시 다른 쪽을 함께 바꿔야 한다.
  \item 시뮬/실기 분기는 전송 대상만 바꾸고 명령 의미는 유지한다.
\end{itemize}
\teachhead{분석할 때 핵심 질문}
\begin{itemize}
  \item fire-and-forget 목록이 양측에서 어긋나면 어떤 증상이 나타날까?
  \item 설정이 \code{system.json}과 \code{localStorage} 두 곳에 있을 때 무엇이 ``정답''인가?
\end{itemize}
\teachhead{점검 요소}
\begin{itemize}
  \item 명령 하나의 형식을 바꿔 보고, 웹\textbf{·}펌웨어 양쪽에서 수정해야 할 지점을 빠짐없이 목록화하라.
\end{itemize}
\end{teachbox}

